% ---
% title: Permutations with repetition
% tags: permutations-with-repetition
% ---

\begin{problem}
	How many different letter arrangements can be made from the letters of
	(a) \mbox{FLUKE} \quad (b) \mbox{PROPOSE} \quad (c) \mbox{MISSISSIPPI} \quad (d) \mbox{ARRANGE}?
\end{problem}

\begin{slides}
	\small
	\exstatement
	\pause
	\sol If all letters were different we would have $n!$ arrangements. Swapping
	identical letters does not change the word, so every word is counted
	$n_1!\,n_2!\cdots$ times: divide by the factorials of the multiplicities.
	\begin{enumerate}
		\item[(a)] 5 different letters: $5! = 120$.
		\item[(b)] P,R,O,P,O,S,E: P and O twice: $\dfrac{7!}{2!\,2!} =
			      \dfrac{5040}{4} = 1260$.
		\item[(c)] M$\times1$, I$\times4$, S$\times4$, P$\times2$: each word is
		      counted $4!\,4!\,2!=1152$ times among the $11!$ orders:
		      $\dfrac{11!}{4!\,4!\,2!} = \dfrac{39\,916\,800}{1152} = 34\,650$.
		\item[(d)] A,R,R,A,N,G,E: A and R twice: $\dfrac{7!}{2!\,2!} = 1260$.
	\end{enumerate}
\end{slides}

\begin{solution}
	\idea If all $n$ letters were different there would be $n!$ arrangements.
	If a letter occurs $m$ times, swapping its copies gives the same word, so
	each word has been counted $m!$ times; we divide by the factorials of all
	multiplicities.

	\step{(a) FLUKE.} 5 different letters: $5! = 120$.

	\step{(b) PROPOSE.} 7 letters: P twice, O twice, R, S, E once.
	\[ \frac{7!}{2!\,2!} = \frac{5040}{4} = 1260. \]

	\step{(c) MISSISSIPPI.} 11 letters: M once, I four times, S four times,
	P twice ($1+4+4+2=11$, always check this). Each word is counted
	$4!\cdot4!\cdot2! = 24\cdot24\cdot2 = 1152$ times among the $11!$ orders:
	\[ \frac{11!}{1!\,4!\,4!\,2!} = \frac{39\,916\,800}{1152} = 34\,650. \]

	\step{(d) ARRANGE.} 7 letters: A twice, R twice, N, G, E once:
	\[ \frac{7!}{2!\,2!}=1260. \]

	\begin{answer}
		(a) $120$ \quad (b) $1260$ \quad (c) $34\,650$ \quad (d) $1260$
	\end{answer}

	\begin{remark}
		\textbf{Common mistake:} dividing by $2!+2!$ instead of $2!\cdot2!$. The
		over-counting factors multiply: the two P's and the two O's can be swapped
		independently.
	\end{remark}
\end{solution}
